\chapter{Fermi Surfaces and Metals}

The Fermi surface is the surface in $\mathbf{k}$-space that separates occupied from unoccupied electron states at $T = 0$. Its shape determines the electronic, thermal, and magnetic properties of metals. In this chapter we present the experimental techniques and data that reveal Fermi surface geometry.

\section{Construction of Fermi Surfaces}

For a free electron gas, the Fermi surface is a sphere of radius $k_F = (3\pi^2 n)^{1/3}$. In real metals, the periodic potential distorts this sphere, especially near Brillouin zone boundaries where band gaps open. The resulting Fermi surface can be remarkably complex.

\section{The de Haas--van Alphen Effect}

The most powerful experimental probe of Fermi surface geometry is the \textbf{de Haas--van Alphen (dHvA) effect}: the oscillation of the magnetic susceptibility (or magnetization) of a metal as a function of $1/B$ at low temperatures.

The Onsager relation connects the oscillation period to the extremal cross-sectional area $A$ of the Fermi surface perpendicular to $\mathbf{B}$:
\begin{equation}
  \Delta\left(\frac{1}{B}\right) = \frac{2\pi e}{\hbar A}
\end{equation}

By rotating the crystal and measuring the oscillation frequency as a function of field direction, the entire Fermi surface can be mapped out.

\subsection{Fermi Surface of Copper}

Copper is the most thoroughly studied metal in Fermi surface physics. Two landmark experiments established its Fermi surface topology:

\paragraph{Pippard (1957).} A.~B.~Pippard~\cite{pippard1957fermi} used the \textbf{anomalous skin effect}---the frequency-dependent surface resistance of a metal---to determine the Fermi surface shape of copper. By measuring the surface impedance as a function of crystallographic orientation, he showed that the Cu Fermi surface is a sphere with ``necks'' protruding toward the $L$ points of the Brillouin zone. This was the first experimental determination of a Fermi surface shape in any metal.

\paragraph{Shoenberg (1962).} D.~Shoenberg~\cite{shoenberg1962dhva} performed comprehensive dHvA measurements on Cu, Ag, and Au single crystals, identifying two principal oscillation frequencies:

\begin{table}[H]
\centering
\caption{De Haas--van Alphen frequencies for copper. ``Belly'' corresponds to the large central cross-section; ``neck'' to the small cross-section at the $L$-point contact with the zone boundary. Data from Shoenberg (1962).}
\label{tab:dhva_cu}
\begin{tabular}{@{}lSS@{}}
\toprule
Orbit & {Frequency ($10^8$ \si{\gauss})} & {$A/A_{\text{free}}$} \\
\midrule
Belly $\langle 111 \rangle$   & 5.90 & 1.00 \\
Belly $\langle 100 \rangle$   & 5.98 & 1.01 \\
Belly $\langle 110 \rangle$   & 5.88 & 1.00 \\
Neck $\langle 111 \rangle$    & 0.22 & 0.037 \\
\bottomrule
\end{tabular}
\end{table}

The belly frequency is nearly isotropic and close to the free-electron value (ratio $\approx 1.00$), confirming that most of the Cu Fermi surface is indeed nearly spherical. The neck frequency is much smaller ($\sim$4\% of belly), corresponding to the narrow tubes connecting the sphere to the zone boundary at the 8 $L$-points.

This topology---a sphere with 8 necks---is one of the most iconic images in solid state physics. It explains:
\begin{itemize}[nosep]
  \item Why Cu is an excellent electrical conductor (large, nearly spherical Fermi surface with high velocity).
  \item Why the optical properties of Cu change abruptly at $\sim$\SI{2}{\electronvolt} (interband transitions from $d$-band to states above $E_F$, giving Cu its characteristic color).
  \item The anomalous positive Hall coefficient observed in some orientations (due to open orbits on the multiply-connected Fermi surface).
\end{itemize}

\section{Experimental Methods Summary}

\begin{table}[H]
\centering
\caption{Experimental methods for Fermi surface studies.}
\label{tab:fs_methods}
\begin{tabular}{@{}lll@{}}
\toprule
Method & What it measures & Resolution \\
\midrule
de Haas--van Alphen  & Extremal cross-section areas & Very high \\
Anomalous skin effect & Fermi surface shape & Moderate \\
ARPES                & $E(\mathbf{k})$ directly & $\sim$\SI{1}{\milli\electronvolt} \\
Positron annihilation & Momentum distribution & Moderate \\
Compton scattering   & Electron momentum density & Low \\
\bottomrule
\end{tabular}
\end{table}

\vspace{1cm}
\noindent\rule{\textwidth}{0.4pt}
\textit{We now turn to superconductivity---a macroscopic quantum phenomenon that was one of the great surprises of condensed matter physics.}
